\documentclass[spanish,letterpaper,oneside]{article}

\def\documenttitle {Ciencias del Aprendizaje: Infografía y Análisis de R. Keith Sawyer (2008)}
\def\documentsubtitle {Fundamentos para la enseñanza y el aprendizaje I · Cohorte SD26}
\def\documentsubject {Actividad 9 · Sesión 3b (Complemento Infográfico Módulo 51440 / 51436)}
\def\documentauthor {Mtro. Martín Jonathan de la Cruz Muñoz}
\def\coursename {Fundamentos para la enseñanza y el aprendizaje I}
\def\coursecode {MAEM-SD26}

\def\universityname {Instituto de Investigación, Innovación y Estudios de Posgrado para la Educación del Estado de Nuevo León}
\def\universityfaculty {Maestría en Aprendizaje y Enseñanza de las Matemáticas}
\def\universitydepartment {Fundamentos para la enseñanza y el aprendizaje I}
\def\universitydepartmentimage {img/iiiepe/marca-iiiepe.png}
\def\universitydepartmentimagecfg {height=1.75cm}
\def\universitylocation {Monterrey, Nuevo León}
\def\studentfullname {Mtro. Martín Jonathan de la Cruz Muñoz}
\def\studentid {00874}
\def\studentemail {martin.delacruz@campus.iiiepe.edu.mx}
\def\programperiod {2do periodo, septiembre--diciembre de 2026 (SD26)}
\def\programmodality {Modalidad presencial}
\def\programshort {MAEM}

\def\authortable {
  \begin{tabular}{ll}
    \textbf{Alumno:} & \studentfullname \\
    Matrícula: & \studentid \\
    Correo institucional: & \studentemail \\
    Unidad: & \coursename \\
    Programa: & \programshort \\
    Periodo: & \programperiod \\
    & \\
    \multicolumn{2}{l}{Fecha de realización: 19 de noviembre de 2026} \\
    \multicolumn{2}{l}{\universitylocation}
  \end{tabular}
}

\input{template}

\usepackage{helvet}
\renewcommand{\familydefault}{\sfdefault}
\usepackage{array}
\usepackage{longtable}
\usepackage{pdflscape}
\usepackage{tikz}
\usepackage{pgfplots}
\pgfplotsset{compat=1.18}
\usepackage[most]{tcolorbox}
\usetikzlibrary{arrows.meta,positioning,calc,fit,backgrounds,shapes.geometric}

\definecolor{iiiepeNavy}{HTML}{163B5C}
\definecolor{iiiepeBlue}{HTML}{1F6F9F}
\definecolor{iiiepeAqua}{HTML}{20B7A8}
\definecolor{iiiepeOrange}{HTML}{F28C28}
\definecolor{iiiepePaper}{HTML}{F6F8F7}
\definecolor{iiiepeInk}{HTML}{1E2A32}
\definecolor{iiiepePale}{HTML}{EEF5F8}
\definecolor{iiiepeMint}{HTML}{E7F6F4}
\definecolor{iiiepeWarm}{HTML}{FFF3E6}
\def\portraittitlecolor {iiiepeNavy}
\setcitestyle{authoryear,open={(},close={)}}
\makeatletter
\renewcommand{\@biblabel}[1]{}
\makeatother
\setlength{\bibhang}{1.27cm}
\setlength{\bibsep}{6pt}
\AtBeginDocument{\renewcommand{\refname}{Referencias}}

\newcommand{\insertcoverwatermark}{%
  \AddToShipoutPictureBG*{%
    \begin{tikzpicture}[remember picture,overlay]
      \node[opacity=0.04,inner sep=0pt] at (current page.center) {%
        \includegraphics[width=0.65\paperwidth]{img/iiiepe/marca-iiiepe.png}%
      };
    \end{tikzpicture}%
  }%
}

\setlength{\parskip}{5pt}\setlength{\parindent}{0pt}
\begin{document}\insertcoverwatermark\templatePortrait\templatePagecfg\templateFinalcfg\fancyfoot{}\fancyfoot[L]{\scriptsize IIIEPE · MAEM · SD26}\fancyfoot[C]{\scriptsize\thepage}\fancyfoot[R]{\scriptsize Actividad 9 · Infografía Sawyer}\setstretch{1.3}\raggedbottom
\section{Introducción y postura central sobre Keith Sawyer (2008)}
La transición hacia una enseñanza matemática comprensiva demanda examinar científicamente cómo aprenden las personas y por qué las estructuras escolares heredadas de la era industrial resultan insuficientes en la sociedad del conocimiento. En el capítulo \textit{Mejorando el aprendizaje: las implicaciones de las investigaciones en las ciencias del aprendizaje}, \citet{sawyer2008optimizing} demuestra que la escuela tradicional —denominada \textit{instruccionismo} o modelo de transmisión y adquisición— se organizó sobre supuestos de sentido común anteriores al desarrollo interdisciplinario de la ciencia cognitiva, la neurociencia educativa, la psicología del desarrollo y los estudios de diseño instruccional.

Mi postura es que superar el modelo estándar en bachillerato técnico (CONALEP y CECyTE) y en formación universitaria en ingeniería no consiste en añadir tecnología superficial ni incrementar la cantidad de ejercicios repetitivos, sino en rediseñar las secuencias matemáticas en torno a comprensión conceptual profunda, activación de saberes previos, articulación externa del razonamiento y andamiaje progresivo \citep{sawyer2008optimizing,star2015algebra}. Por tanto, este informe integra una infografía panorámica de las tesis de Sawyer seguida de su análisis crítico y transferencia concreta a la modelación matemática técnica.

\clearpage
\begin{landscape}
\thispagestyle{empty}
\section*{Infografía · Ciencias del Aprendizaje (R. Keith Sawyer, 2008)}
\vspace{-4pt}
\begin{center}
\resizebox{0.97\linewidth}{!}{%
\begin{tikzpicture}[font=\scriptsize, >=Latex]
  % Marco principal apaisado exacto (-11.4, -7.05) a (11.4, 7.25)
  \fill[iiiepePaper] (-11.4,-7.05) rectangle (11.4,7.25);
  \draw[iiiepeNavy,very thick,rounded corners=6pt] (-11.4,-7.05) rectangle (11.4,7.25);

  % Encabezado
  \fill[iiiepeNavy,rounded corners=5pt] (-11.3,5.48) rectangle (11.3,7.15);
  \node[anchor=west,text width=17.4cm,align=left] at (-11.0,6.32) {%
    {\color{iiiepeAqua!45!white}\scriptsize\bfseries INFOGRAFÍA ACADÉMICA · R. KEITH SAWYER (2008) · MÓDULO 51436 / 51440 (SD26)\par}%
    \vspace{2pt}%
    {\color{white}\normalsize\bfseries MEJORANDO EL APRENDIZAJE: IMPLICACIONES DE LAS CIENCIAS DEL APRENDIZAJE\par}%
    \vspace{2pt}%
    {\color{white!85}\scriptsize Del instruccionismo de la era industrial hacia ambientes de comprensión profunda, andamiaje y conocimiento situado\par}
  };
  \node[anchor=east] at (10.95,6.32) {\includegraphics[height=1.05cm]{img/iiiepe/logo-iiiepe-blanco.png}};

  % BLOQUE 1: Instruccionismo vs Economia del Conocimiento
  \node[draw=iiiepeNavy,fill=white,thick,rounded corners=4pt,text width=6.55cm,minimum height=4.42cm,anchor=north west,inner sep=5pt,align=left,font=\fontsize{6.65}{8.15}\selectfont] (b1) at (-11.05,5.25) {%
    {\colorbox{iiiepeNavy}{\color{white}\bfseries\ 1. CRISIS DEL MODELO ESTÁNDAR (INSTRUCCIONISMO) \ }}\\[4pt]
    \textbf{· Supuestos no probados:} El conocimiento es una colección de hechos y procedimientos paso a paso; el docente transmite y el examen mide recuerdo memorístico.\\[3pt]
    \textbf{· Modelo fabril (Papert):} La secuencia «simple a complejo» fue definida por adultos expertos, sin investigar cómo razonan realmente niños y jóvenes.\\[3pt]
    \textbf{· Economía creativa y conocimiento situado:} Los «trabajadores del conocimiento» (analistas simbólicos) crean e improvisan colaborando con herramientas y prácticas reales.
  };

  % BLOQUE 2: Hallazgos clave
  \node[draw=iiiepeBlue,fill=white,thick,rounded corners=4pt,text width=6.85cm,minimum height=4.42cm,anchor=north,inner sep=5pt,align=left,font=\fontsize{6.6}{8.1}\selectfont] (b2) at (0.0,5.25) {%
    {\colorbox{iiiepeBlue}{\color{white}\bfseries\ 2. HALLAZGOS DE LAS CIENCIAS DEL APRENDIZAJE \ }}\\[4pt]
    \textbf{· Comprensión profunda vs. superficial:} Hechos y algoritmos solo son útiles al integrarse en estructuras conceptuales transferibles a situaciones nuevas.\\[3pt]
    \textbf{· Expertis cognitiva:} Combinar esquemas mentales ricos, planes flexibles, improvisación contextual y monitoreo del propio pensamiento.\\[3pt]
    \textbf{· Saberes previos («matemáticas ingenuas»):} El aula no recibe mentes vacías; los errores nacen de teorías previas que deben hacerse visibles para reestructurarse.
  };

  % BLOQUE 3: Articulacion, Reflexion y Andamiaje
  \node[draw=iiiepeAqua!80!black,fill=iiiepeMint,thick,rounded corners=4pt,text width=6.55cm,minimum height=4.42cm,anchor=north east,inner sep=5pt,align=left,font=\fontsize{6.65}{8.15}\selectfont] (b3) at (11.05,5.25) {%
    {\colorbox{iiiepeAqua!80!black}{\color{white}\bfseries\ 3. ARTICULACIÓN, METACOGNICIÓN Y ANDAMIAJE \ }}\\[4pt]
    \textbf{· Externalizar para aprender:} Pensar en voz alta, escribir argumentos y explicar hipótesis construye comprensión más profunda que escuchar en silencio.\\[3pt]
    \textbf{· Reflexión metacognitiva:} Examinar el propio camino seguido permite ajustar heurísticas dentro del espacio del problema.\\[3pt]
    \textbf{· Andamiaje (\textit{scaffolding}):} Ayuda personalizada mediante pistas, modelado o preguntas que sostiene la obra y se retira gradualmente hasta lograr autonomía.
  };

  % BLOQUE 4: Principios de diseno de ambientes
  \node[draw=iiiepeOrange!90!black,fill=iiiepeWarm,thick,rounded corners=4pt,text width=10.55cm,minimum height=4.05cm,anchor=south west,inner sep=6pt,align=left,font=\fontsize{6.75}{8.25}\selectfont] (b4) at (-11.05,-5.82) {%
    {\colorbox{iiiepeOrange!85!black}{\color{white}\bfseries\ 4. PRINCIPIOS DE DISEÑO Y CURRÍCULO PROFUNDO (LECCIÓN TIMSS) \ }}\\[4pt]
    \textbf{· Personalización y fuentes diversas:} Respetar ritmos y estructuras cognitivas previas apoyándose en libros, simulaciones, expertos y tecnología bien diseñada; el profesor deja de ser única fuente.\\[3pt]
    \textbf{· Conocimiento distribuido y colaborativo:} En la vida profesional la inteligencia opera en equipos y artefactos compartidos, acelerando la construcción conceptual frente al aislamiento escolar.\\[3pt]
    \textbf{· Coherencia frente a saturación temática:} Estudios como TIMSS muestran que currículos saturados de temas superficiales logran menor desempeño que enfoques coherentes centrados en comprender.
  };

  % BLOQUE 5: Rol docente, evaluacion y transferencia
  \node[draw=iiiepeNavy,fill=iiiepePale,thick,rounded corners=4pt,text width=10.55cm,minimum height=4.05cm,anchor=south east,inner sep=6pt,align=left,font=\fontsize{6.75}{8.25}\selectfont] (b5) at (11.05,-5.82) {%
    {\colorbox{iiiepeNavy}{\color{white}\bfseries\ 5. EVALUACIÓN AUTÉNTICA, TUTELADO Y ROL DOCENTE \ }}\\[4pt]
    \textbf{· Superar pruebas descontextualizadas:} Las evaluaciones estandarizadas miden memoria superficial e ignoran errores conceptuales previos; urge evaluar integración y resolución.\\[3pt]
    \textbf{· Docente trabajador del conocimiento:} Requiere dominio teórico-disciplinar, autonomía profesional y conducción de taller experto en lugar de control autocrático de fábrica.\\[3pt]
    \textbf{· Aprendizaje tutelado (\textit{apprenticeship}) y uso en Matemáticas:} Aprender junto a prácticas reales de modelación: tablas, gráficas, unidades y verificación crítica en ingeniería y bachillerato técnico.
  };

  % Conectores limpios
  \draw[->,very thick,iiiepeNavy] (b1.east) -- (b2.west);
  \draw[->,very thick,iiiepeBlue] (b2.east) -- (b3.west);
  \draw[->,very thick,iiiepeOrange!85!black] (b1.south) -- (b4.north west);
  \draw[->,very thick,iiiepeAqua!80!black] (b3.south) -- (b5.north east);
  \draw[<->,very thick,iiiepeNavy] (b4.east) -- (b5.west);

  % Pie de fuentes y creditos
  \fill[iiiepeNavy,rounded corners=4pt] (-11.3,-6.95) rectangle (11.3,-6.02);
  \node[anchor=west,text width=14.0cm,align=left,font=\tiny\color{white}] at (-11.0,-6.48) {%
    \textbf{FUENTE BASE (APA 7):} Sawyer, R. K. (2008). Optimizing learning: Implications of learning sciences research. En OECD/CERI, \textit{Innovating to learn, learning to innovate} (cap. 2, pp. 45--65). OECD Publishing. \url{https://doi.org/10.1787/9789264039958-en}
  };
  \node[anchor=east,text width=7.5cm,align=right,font=\tiny\color{white}] at (11.0,-6.48) {%
    \textbf{CRÉDITOS:} Mtro. Martín Jonathan de la Cruz Muñoz (Matrícula 00874)\\
    MAEM · IIIEPE · Cohorte SD26 · 19 de noviembre de 2026
  };
\end{tikzpicture}%
}
\end{center}
\end{landscape}
\clearpage

\section{Análisis crítico del capítulo de Keith Sawyer (2008)}
La aportación medular de \citet{sawyer2008optimizing} consiste en demostrar que el \textit{modelo estándar} de escolarización responde históricamente a las necesidades de estandarización de la fábrica del siglo XIX y principios del XX, pero fracasa cuando la sociedad exige trabajadores capaces de trabajar creativamente con conceptos complejos. A partir del texto de la Sesión 3, destacan cuatro implicaciones estructurales para nuestra enseñanza matemática:

\begin{enumerate}
  \item \textbf{Las matemáticas ingenuas y las concepciones previas no son ignorancia simple:} \citet{sawyer2008optimizing} retoma la premisa piagetiana y constructivista de que los aprendices poseen estructuras cualitativamente distintas de conocimiento previo. En nuestros planteles técnicos y de ingeniería, cuando un estudiante confunde la ordenada al origen $b$ con la pendiente $m$ o aplica proporcionalidad directa $y=kx$ a fenómenos con sesgo inicial $y=mx+b$, no basta con repetir la definición correcta; es indispensable diseñar situaciones contrastantes que vuelvan explícito y discutible el error conceptual.
  \item \textbf{Articulación y andamiaje efectivo (\textit{scaffolding}):} Comprender profundamente requiere que el alumnado exteriorice su pensamiento en desarrollo mediante diagramas, justificaciones orales y discusiones escritas \citep{sawyer2008optimizing}. Sin embargo, el andamiaje auténtico nunca consiste en resolver el problema por el estudiante, sino en ofrecer pistas graduadas —tablas parcialmente estructuradas, preguntas sobre unidades dimensionales, contraejemplos— que se retiran conforme emerge la autonomía.
  \item \textbf{Coherencia frente a saturación curricular:} Al analizar los resultados comparativos internacionales (TIMSS), \citet{sawyer2008optimizing} evidencia que los programas que acumulan listas extensas de temas inconexos generan menor logro matemático que aquellos orientados hacia ideas nucleares profundas e interconectadas. Dado que en bachillerato técnico el tiempo de aula es acotado, priorizar pocas familias de modelos estudiadas desde múltiples representaciones resulta más fructífero que cubrir mecánicamente decenas de fórmulas aisladas.
  \item \textbf{Evaluación de comprensión e integración frente al recuerdo superficial:} Las pruebas convencionales centradas en reproducir algoritmos descontextualizados ocultan si el estudiante puede elegir entre estrategias equivalentes o detectar el límite físico de validez de una ecuación \citep{sawyer2008optimizing,star2015algebra}. Por ello, evaluar desde las ciencias del aprendizaje implica valorar la coherencia entre tabla, gráfica, ecuación y argumentación escrita durante el proceso mismo \citep{eef2021feedback}.
\end{enumerate}

\section{Contrastación en la enseñanza matemática técnica e ingeniería}
En mi práctica con grupos de CONALEP, CECyTE e ingeniería, los hallazgos de \citet{sawyer2008optimizing} se traducen en decisiones didácticas precisas comparadas en la tabla siguiente:

{\small\setstretch{1.15}\renewcommand{\arraystretch}{1.32}
\begin{longtable}{|>{\raggedright\arraybackslash}p{0.22\linewidth}|>{\raggedright\arraybackslash}p{0.35\linewidth}|>{\raggedright\arraybackslash}p{0.37\linewidth}|}\hline
\rowcolor{iiiepePale} \textbf{Dimensión didáctica} & \textbf{Instruccionismo tradicional} & \textbf{Ciencias del aprendizaje (Sawyer, 2008)} \\\hline
\endfirsthead
\hline
\rowcolor{iiiepePale} \textbf{Dimensión didáctica} & \textbf{Instruccionismo tradicional} & \textbf{Ciencias del aprendizaje (Sawyer, 2008)} \\\hline
\endhead
Rol del conocimiento previo & Se asume mente vacía o repaso mecánico de fórmulas previas. & Diagnóstico explícito de concepciones erróneas (\textit{naive math}) antes de formalizar \citep{sawyer2008optimizing}. \\\hline
Gestión de problemas matemáticos & Ejercicios repetitivos tipo receta («despeje $x$ en veinte ecuaciones»). & Modelación situada con verificación dimensional e interpretación de parámetros \citep{conalep2025mecatronica}. \\\hline
Apoyo docente (\textit{scaffolding}) & El profesor muestra la solución completa cuando el grupo se bloquea. & Pistas graduadas que retienen el desafío cognitivo y se retiran progresivamente \citep{eef2017mathematics}. \\\hline
Colaboración y fuentes & Trabajo individual silencioso con libro único y apuntes dictados. & Conocimiento distribuido: tríos con roles rotativos, hoja de datos y contraste gráfico \citep{sawyer2008optimizing}. \\\hline
Evaluación del aprendizaje & Examen final acumulativo centrado en puntaje de respuestas cerradas. & Retroalimentación procesual sobre justificación escrita y segunda versión corregida \citep{eef2021feedback}. \\\hline
\end{longtable}}

En un problema ilustrativo por realizar sobre calibración de un sensor industrial ($y = 1.05x + 0.2$), el aprendizaje superficial se limita a evaluar $y(40) = 42.2$; en contraste, una secuencia basada en las ciencias del aprendizaje solicita predecir primero si el sesgo relativo crece o decrece con la magnitud medida, articular por escrito la diferencia entre error absoluto ($0.05x + 0.2$) y error porcentual relativo, discutir en equipo el rango tolerable de operación y justificar qué corrección lineal requeriría el instrumento. Asimismo, estas pautas coinciden con la democratización escolar propuesta por \citet{meirieu2020escuela} y con el perfil humanista del currículo nacional \citep{sep2023pensamiento}.

\clearpage
\section{Conclusiones}
En conclusión, sostengo que la obra de \citet{sawyer2008optimizing} provee el cimiento empírico y cognitivo indispensable para transformar el aula de matemáticas técnicas sin caer en modas pedagógicas carentes de sustento. Considero que articular comprensión conceptual profunda, exteriorización argumentada del razonamiento, trabajo colaborativo distribuido y evaluación formativa continua permite superar la fragmentación del instruccionismo fabril.

Por ello, daría prioridad en mis próximos bloques de enseñanza a diseñar menos rutinas mecánicas repetitivas y más situaciones problemáticas auténticas con andamiaje transparente, verificando mediante producciones escritas y discusiones grupales que cada estudiante logre transferir los modelos lineales hacia la toma fundamentada de decisiones profesionales y ciudadanas.

\clearpage
\begin{thebibliography}{99}
\bibitem[Collin y Quigley(2021)]{eef2021feedback} Collin, J., y Quigley, A. (2021). \textit{Teacher feedback to improve pupil learning}. Education Endowment Foundation. \url{https://educationendowmentfoundation.org.uk/education-evidence/guidance-reports/feedback}
\bibitem[CONALEP(2025)]{conalep2025mecatronica} Colegio Nacional de Educación Profesional Técnica. (2025). \textit{Perfil de egreso: Profesional Técnico-Bachiller en Mecatrónica. Modelo Académico 2025}. Secretaría Académica, CONALEP.
\bibitem[Henderson et~al.(2017)]{eef2017mathematics} Henderson, P., Hodgen, J., Foster, C., y Kuchemann, D. (2017). \textit{Improving mathematics in Key Stages 2 and 3}. Education Endowment Foundation.
\bibitem[Meirieu(2020)]{meirieu2020escuela} Meirieu, P. (2020, 18 de abril). ``La escuela después''... ¿Con la pedagogía de antes? \textit{Movimiento Cooperativo de Escuela Popular}. \url{https://www.mcep.es/la-escuela-despues-con-la-pedagogia-de-antes-philippe-meirieu/}
\bibitem[Sawyer(2008)]{sawyer2008optimizing} Sawyer, R. K. (2008). Optimizing learning: Implications of learning sciences research. En OECD, \textit{Innovating to learn, learning to innovate} (pp. 45--65). OECD Publishing. \url{https://doi.org/10.1787/9789264039958-en}
\bibitem[SEP(2023)]{sep2023pensamiento} Secretaría de Educación Pública. (2023). \textit{Programa de estudios del recurso sociocognitivo Pensamiento Matemático I}. Subsecretaría de Educación Media Superior, SEP.
\bibitem[Star et~al.(2015)]{star2015algebra} Star, J. R., Caronongan, P., Foegen, A., Furgeson, J., Keating, B., Larson, M. R., Lyskawa, J., McCallum, W. G., Porath, J., y Zbiek, R. M. (2015). \textit{Teaching strategies for improving algebra knowledge in middle and high school students} (NCEE 2015-4010). Institute of Education Sciences.
\end{thebibliography}

\end{document}
